\section{语言模型：从词向量到后训练}

前两章解释了模型范式和训练底座，现在回到今天大多数人最熟悉的入口：语言模型。本章把“表示、预训练、后训练”连成一条线。语言模型的发展可以看成三个阶段：先把词变成向量，再用预训练学习语言规律，最后用指令、偏好和开源后训练把模型变得可用、可控、可服务用户。

\begin{figure}[H]
\centering
\includegraphics[width=0.96\textwidth]{language-model-lineage.png}
\caption{语言模型谱系：从词向量到预训练，再到指令对齐与开源后训练。自制概念图，依据 02:21:29--03:08:08 对谈内容整理。}
\end{figure}

\begin{knowledgebox}{读图：语言模型的主线是“表示到行为”}
Word2Vec 解决词的表示，神经机器翻译证明深度模型能在线服务，GPT/BERT 建立预训练范式，GPT-3 展示规模涌现，InstructGPT 和 Tulu 3 则把模型推向人类可用和开源后训练。
\end{knowledgebox}

\subsection{Word2Vec 与 Google Translate：表示和部署}

本章先从语言模型的前史讲起，因为没有表示学习，就没有后来的上下文学习。Word2Vec 用机器学习把词转成向量，让语义关系进入连续空间。Google Translate 的神经网络部署则展示了深度学习不只在论文里有效，也能在大规模线上系统中替代传统方法。这两步分别回答了“语言怎样表示”和“神经模型怎样服务真实用户”。

\begin{importantbox}{词向量的意义}
词向量把离散词映射到连续空间，使相似词在向量空间中更接近。这为后来的上下文表示、Transformer embedding 和语义检索铺路，但它仍然是静态表示，无法像 LLM 那样按上下文动态理解词义。
\end{importantbox}

\subsection{GPT-1、BERT、GPT-2 与 GPT-3}

从词向量和线上翻译往后走，语言模型真正进入预训练时代。本节的 GPT-1、BERT、GPT-2 和 GPT-3 分别代表不同阶段：GPT-1 提出无监督预训练加监督微调的 NLP 新范式；BERT 用双向 masked language modeling 在理解任务上成为“曾经的王”；GPT-2 强化了生成式预训练的路线，并让大家开始认真看待“告别微调”的可能；GPT-3 则把 Scaling Law、组织押注和大规模工程结合起来，成为 ChatGPT 前夜最重要的信号。

\begin{knowledgebox}{GPT 与 BERT 的分岔}
\begin{tabularx}{\textwidth}{p{0.22\textwidth}p{0.34\textwidth}X}
\toprule
路线 & 训练方式 & 典型优势 \\
\midrule
BERT & 遮盖词预测，双向上下文 & 理解、分类、抽取等判别任务强。 \\
GPT & 自回归 next-token prediction & 生成、补全、对话和通用任务扩展性强。 \\
GPT-3 & 大规模自回归预训练 & 少样本、上下文学习和涌现能力成为焦点。 \\
\bottomrule
\end{tabularx}
\end{knowledgebox}

\subsection{InstructGPT 与 Tulu 3：让模型进入文明社会}

InstructGPT 的核心意义是把模型从“会续写互联网文本”推向“能按人类意图回答”。这背后是指令数据、偏好数据和 RLHF 等后训练方法。Tulu 3 则代表开源社区试图把后训练流程系统化、透明化，让模型能力不只掌握在少数闭源实验室手中。

\begin{warningbox}{后训练不是锦上添花}
预训练给模型世界知识和语言能力，后训练决定它如何响应人类、遵守指令、拒绝危险请求、按格式输出和完成任务。没有后训练，强模型也可能很难变成好产品。
\end{warningbox}

\subsection{本章小结}

语言模型线把 AI 从表示学习推到通用交互入口。Word2Vec 到 GPT-3 是能力增长，InstructGPT 到 Tulu 3 是可用性、对齐和开源后训练的增长。今天的 LLM 产品建立在这两条增长曲线交汇处。

